\section{Introduction}
Cooperative trajectory planning enables multiple unmanned aerial vehicles (UAVs) to execute coordinated missions in a shared three-dimensional workspace. When several vehicles approach an intersection or a narrow passage, their trajectories must satisfy obstacle clearance, mutual separation, and motion limits simultaneously. Aerial-swarm research has established the importance of these coupled requirements~\cite{chung2018}, while agile quadrotor teams and coordinated formation transitions demonstrate their practical relevance~\cite{kushleyev2012,preiss2017}. Recent work combines three-dimensional reciprocal velocity obstacles with multi-agent reinforcement learning to coordinate collision avoidance within structured urban airspace~\cite{zhong2025rvo}.

Computational efficiency is essential when a feasible coordinated motion must be obtained within a limited planning interval. Enhanced sequential convex programming has been integrated with model predictive control to address formation tracking and correct linearization errors that can compromise feasibility~\cite{liu2025escp}. Conflict-free multi-robot planning coupled with predictive control connects coordination decisions to feasible motion~\cite{wu2025conflict}. Cooperative exploration combines dynamic Voronoi partitioning with learned obstacle avoidance~\cite{zhao2025exploration}, while multi-UAV coverage planning jointly considers task allocation and travel time under repeated round-trip operation~\cite{li2025coverage}. Together, these studies motivate methods that account for coupled coordination, safety, and timing requirements while keeping computation manageable.

High-conflict regions arise where assigned routes intersect or several UAVs compete for a shared passage around obstacles. Safe corridors express collision avoidance through conservative linear restrictions. A useful reference can nevertheless produce local corridors that conflict with one another or with bounded motion. Under a prescribed completion time, all of these requirements must be reconciled within the original flight horizon. A failed local solve therefore calls for a mechanism that can revise the coordinated motion and its surrounding geometry while preserving the task boundaries.

This paper develops a lightweight iterative optimization method (LIOM) for fixed-horizon, rest-to-rest cooperative flight. The double-integrator model and fixed spatial and relative corridors give a convex quadratic structure. Separate penalties on kinematic defects and relative-box violations allow intermediate candidates to reshape the corridors while retaining hard static free-space and translational motion bounds. A synchronized reference initializes the process. Bounded penalty updates support numerical recovery, and a closed-form endpoint control correction connects recovered candidates to continuous-velocity motion over the prescribed horizon.

The main contributions are twofold:
\begin{enumerate}
\item A fixed-horizon recovery formulation reconciles prescribed rest-to-rest states with spatial and relative collision-avoidance regions. The affine model and corridor constraints retain a convex quadratic subproblem, whose nonemptiness is established by an explicit feasible-point construction under endpoint-compatible spatial boxes. This formulation permits continued local refinement through temporary incompatibilities among kinematics, static avoidance, and mutual separation.
\item A closed-form endpoint control correction connects numerical recovery to continuous-velocity motion with the prescribed boundary states. Exact integration is followed by complete-interval clearance and motion-limit checks on the reconstructed trajectory; these checks determine final acceptance. The procedure retains the original flight horizon and requires no additional iterative optimization solve for endpoint correction.
\end{enumerate}
Simulation studies evaluate accepted cooperative motions, computation time, and the role of initialization. Task-matched comparisons and controlled pressure sweeps examine how local geometric restrictions and recovery affect success within a fixed flight horizon.
